% Measurement against the transistor-level simulation of the ADC. Macros from results_sim.tex
% (written by make_report.py from the sim_summary.json of saradc_auto/sim/report/make_sim_report.py).
\input{results_sim.tex}
\section{Comparison with simulation}
\label{sec:comparison}
The ADC was simulated at transistor level in \texttt{saradc\_auto} (data set\\
\texttt{\SimSource}):
ngspice-47~\cite{ngspice} with the PSP103 models of the PDK~\cite{ihppdk} (\SimPdk), typical corner,
\SI{25}{\celsius}, no noise and no mismatch; the post-route schematic netlist\ifdefined\PexBitEA{} and, for the
static transfer, the post-layout netlist of the taped-out layout (Section~\ref{sec:cmp_pex})\fi.
The operating point is that of the adapter board: supply \SI{1.2}{V}, references
\SI{\BoardVrefn}{V}/\SI{\BoardVrefp}{V}, input common mode \SI{0.6}{V}, clock \SI{10}{MHz} with 7
clocks per conversion. The simulated input is the differential voltage at the pins of the macro, so its
LSB is twice the LSB at the input of the board ($v_d = 2(v_{in} - \SI{0.6}{V})$); both are compared in
their own end-point LSB below.

\ifdefined\PexBitEA
The layout parasitics come from a split extraction of the taped-out GDS (details in the simulation
report): neither Magic~\cite{magic} (capacitances two orders of magnitude too large, no comparison
possible in Netgen~\cite{netgen}) nor KLayout-PEX~\cite{kpex} on the whole layout (single-threaded
2.5D engine, runs of up to \SI{16}{h} without result) gave a usable netlist, but KLayout-PEX 0.3.12
does extract small windows. The layout is cut into 81 windows along fillers and taps, every window is
extracted (2.5D, coupling capacitances, no resistance) and the windows are joined again by net name.
The joined netlist has the 24\,198 transistors of the schematic and 68\,519 capacitors, and matches
the schematic in the transistor widths on every one of its 572 named nets.
\else
The layout parasitics are not included because neither extractor gave a usable netlist (details in the
simulation report):
\begin{itemize}
  \item \textbf{Magic/Netgen}~\cite{magic,netgen}: magic 8.3.623 extracts the same devices as the
        schematic, but its capacitances are two orders of magnitude too large (\SI{4.3}{nF} on AVDD,
        about \SI{415}{pF} on the output of a single switch cell, against about \SI{10}{pF} for the whole
        DAC), and Netgen 1.5.287 cannot compare it with the schematic because the transistors are
        written as calls to undefined \texttt{sg13\_lv\_nmos}/\texttt{pmos} subcircuits.
  \item \textbf{KLayout-PEX}~\cite{kpex} 0.3.12: its LVS needs a schematic with the standard cells
        flattened (their layouts have no pin labels), and its 2.5D engine, single-threaded on the flat
        layout, did not finish in runs of \SI{4}{h}, \SI{12}{h} and \SI{10}{h}, nor in a fourth run (fillers
        removed) stopped after \SI{16}{h}.
\end{itemize}
The post-layout comparison will be added when an extraction is available.
\fi

\subsection{Static transfer}
\begin{figure}[H]\centering
\includegraphics{fig/cmp_dnl_inl.pdf}
\caption{DNL and INL of the chip (ramp test) and of the simulation (bisection of the transition
levels) of the schematic\ifdefined\PexBitEA{} and of the post-layout netlist\fi, each in its own end-point LSB.}
\label{fig:cmp_static}
\end{figure}
The schematic converter is ideal within the resolution of the simulation: DNL from \num{\SimDnlMin} to
\num{\SimDnlMax} LSB, INL from \num{\SimInlMin} to \num{\SimInlMax} LSB (\SimInlRms{} LSB rms), no missing codes,
binary weights within \SI{0.1}{\percent}. The chip has the same functionality (monotonic, no missing
codes) but not the same linearity:
\begin{center}\small
\begin{tabular}{lccccc}
\toprule
Bit & 0 & 1 & 2 & 3 & 4 \\
\midrule
Weight, measured (LSB) & \BitWA & \BitWB & \BitWC & \BitWD & \BitWE \\
Error, measured (\%) & \BitEA & \BitEB & \BitEC & \BitED & \BitEE \\
Weight, simulated (LSB) & \SimBitWA & \SimBitWB & \SimBitWC & \SimBitWD & \SimBitWE \\
Error, simulated (\%) & \SimBitEA & \SimBitEB & \SimBitEC & \SimBitED & \SimBitEE \\
\ifdefined\PexBitEA
Weight, post-layout (LSB) & \PexBitWA & \PexBitWB & \PexBitWC & \PexBitWD & \PexBitWE \\
Error, post-layout (\%) & \PexBitEA & \PexBitEB & \PexBitEC & \PexBitED & \PexBitEE \\
\fi
\bottomrule
\end{tabular}
\end{center}
The measured INL (\InlRms{} LSB rms) differs from the schematic one by \CmpInlRms{} LSB rms, i.e.\ all of
it.\ifdefined\PexBitEA{} The post-layout netlist does not: its weight errors have the sign of the measured
ones on all five bits, its INL is \PexInlRms{} LSB rms, and it differs from the measured INL by
\PexCmpInlRms{} LSB rms (correlation \PexInlCorr{} for the INL, \PexDnlCorr{} for the DNL);
Section~\ref{sec:cmp_pex} explains where the errors come from.\fi{} The gain error is \SI{\FeGainErr}{\percent} measured against \SI{\SimGainErr}{\percent} simulated
(schematic\ifdefined\PexBitEA; \SI{\PexGainErr}{\percent} post-layout, where the wiring capacitance of the
top plates attenuates the DAC\fi);
the measured value also contains the tolerance of the front-end gain (R3/R4) and of the reference
voltages, which are known only to the resolution of the multimeter (\SI{10}{mV} on either is
\SI{2.5}{\percent} of the span), so it does not by itself show a difference of the chip. The offset of
\num{\FeOffsetLsb} LSB is not in the simulation either (\num{0.00} LSB); it comes from the board (the
\SI{0.6}{V} divider and the op-amp offsets, doubled by the front end) or from the offset of the comparator,
which a simulation without mismatch cannot show.

\ifdefined\PexBitEA
\subsection{Post-layout simulation: where the weight errors come from}
\label{sec:cmp_pex}
Each bit of the DAC is a group of MOS capacitors between a switched plate (VSH, connected to the input,
to VOUTH or to VOUTL) and a floating plate (FL); VOUTH and VOUTL are the inputs of the comparator. The
unit capacitor is \SI{67.8}{fF} with FL floating. The wiring of the plates adds capacitance of the same
order (Table~\ref{tab:par}): about \SI{100}{fF} per unit on VSH and 65 to \SI{75}{fF} on FL, mostly to
the supplies, which only scales the weights, but also several femtofarads per unit directly to VOUTH and
VOUTL, which changes how much charge each bit moves onto the comparator inputs, and that part is
different for every bit:
\begin{itemize}
  \item bits 1 and 3 have extra bridges on Metal5/TopMetal1 and about \SI{10}{\percent} more wiring capacitance per
        unit on both plates;
  \item the switched plate of the MSB couples \SI{11.3}{fF} per unit to the comparator input of its
        own half, the other bits about \SI{5}{fF} to each input;
  \item the floating plates of bits 0 to 3 couple 7 to \SI{10}{fF} per unit to the comparator input of
        the other half, those of the MSB less than \SI{1}{fF};
  \item the two halves are not symmetric: the floating plates of the upper half couple about \SI{10}{fF}
        to VOUTH and less than \SI{1}{fF} to VOUTL, those of the lower half 3 and \SI{7}{fF} to VOUTH and VOUTL.
\end{itemize}
With all parasitics the taped-out netlist gives weight errors of
\PexBitEA/\PexBitEB/\PexBitEC/\PexBitED/\PexBitEE\,\% (bit 0 to 4, measured
\BitEA/\BitEB/\BitEC/\BitED/\BitEE\,\%).
\ifdefined\PexPlBitEA
With only the capacitors on the plates it gives \PexPlBitEA/\PexPlBitEB/\PexPlBitEC/\PexPlBitED/\PexPlBitEE\,\%%
\ifdefined\PexNpBitEA{} and with all the other capacitors (top plates, switches, controls, without the plates)
\PexNpBitEA/\PexNpBitEB/\PexNpBitEC/\PexNpBitED/\PexNpBitEE\,\%\fi: each part alone already has the
pattern of the measurement (the two do not simply add: their sum is larger than the errors of the full
netlist). The plate wiring also
makes the attenuation of the DAC (the gain error); the other part is the capacitance of the top
plates and of the switch control nets, which couple the switching of each bit into the comparator
inputs.
\fi
The same circuit placed again with one half of the DAC as a macro mirrored into the other half
(hierarchical re-layout, not fabricated) gives \HierBitEA/\HierBitEB/\HierBitEC/\HierBitED/\HierBitEE\,\%
and \HierInlRms{} LSB rms of INL instead of \PexInlRms{}: the symmetry of the two halves halves the
INL, but the bit-to-bit differences of the plate wiring remain.
\IfFileExists{table_par.tex}{%
\begin{table}[H]\centering\scriptsize
\input{table_par.tex}
\caption{Wiring capacitance of the plates of the taped-out layout per unit capacitor (\si{fF}), split
by the net on the other side. H is the lower half (MSB on VOUTH), L the upper half (MSB on VOUTL, the
other bits with VOUTH and VOUTL exchanged). VSH--FL is the coupling between the two plates of the same
bit, in parallel with the unit capacitors.}
\label{tab:par}
\end{table}}{}
\fi

\subsection{Effective number of bits}
The ENOB measured at low frequency, \SineRefEnob{} bits at \SI{\SineRefFreq}{kHz}, is what the measured
static transfer predicts (Section~\ref{sec:board}). Quantizing the same fitted input with the simulated
transition levels, scaled to the measured end points, gives \SimStSinad{}~dB, i.e.\ \SimStEnob{} bits,
practically the \SI{30.2}{dB} of an ideal \AdcBits-bit quantizer at this amplitude. The schematic
therefore allows about one bit more than the chip delivers, and the whole difference is the static
nonlinearity, the binary weight errors of the DAC.\ifdefined\PexBitEA{} With the post-layout transition
levels the same input gives \PexStSinad~dB, \PexStEnob{} bits: the layout parasitics account for about
half of the lost bit.\fi

\IfFileExists{table_cmp.tex}{%
\subsection{Dynamic performance}
The sines were simulated in two ways: with the ideal differential input at the pins of the macro, and
with each input given the amplitude and phase of the linear model of the front end
(Section~\ref{sec:board}, GBW \SI{\FeGbw}{MHz}). The second is the one to compare with the measurement;
it contains the amplitude loss and the common-mode ripple of the board, but not the slew rate or other
nonlinearity of the op-amps.
\begin{figure}[H]\centering
\includegraphics{fig/cmp_enob.pdf}
\caption{ENOB of the chip and of the transient simulation, with the ideal differential input and with
the inputs of the front-end model.}
\label{fig:cmp_enob}
\end{figure}
\begin{figure}[H]\centering
\includegraphics{fig/cmp_harmonics.pdf}
\caption{Second and third harmonic, chip and simulation. With the ideal input H2 cancels exactly
(below \SI{-100}{dBc}) and is not drawn.}
\label{fig:cmp_harm}
\end{figure}
\begin{table}[H]\centering\small
\input{table_cmp.tex}
\caption{Dynamic results, measured and simulated.}
\label{tab:cmp}
\end{table}
\begin{itemize}
  \item \textbf{The converter has no dynamic loss.} With the ideal input the simulated ENOB goes from
        \SimEnobLow{} to \SimEnobHigh{} bits at \SI{\SimFreqHigh}{kHz}: the sampling switches and the
        MOS capacitors do not limit the converter up to the Nyquist frequency.
  \item \textbf{The drop of ENOB with frequency is the board.} With the front-end model the simulated
        ENOB falls to \SimFeEnobHigh{} bits at \SI{\SimFreqHigh}{kHz}, mostly because the amplitude
        falls to \SimFeAmpHigh{} LSB. That is a loss of about \num{0.7} bits, the same as the measured
        loss from \SineRefEnob{} to \MeasEnobHigh{} bits at \SI{\MeasFreqHigh}{kHz}. The measured curve
        lies about one bit below the simulated one at every frequency, which is the static
        nonlinearity of the previous subsection.
  \item \textbf{The measured H2 is distortion of the front end.} The common-mode ripple of the linear
        front-end model produces at most \SI{\SimFeHtwoMax}{dBc} of H2 in the simulated converter. The
        measurement reaches \SI{\MeasHtwoHigh}{dBc} at \SI{\MeasFreqHigh}{kHz}, so the ADC's
        response to the common-mode ripple cannot explain it: the op-amps themselves distort near
        their slew-rate limit, which also explains why the measured amplitude at high frequency is
        below the linear model (Section~\ref{sec:board}).
  \item \textbf{H3 and H5 are static.} The simulated H3 stays between \SI{-63}{dBc} and
        \SI{-43}{dBc}; the measured H3 and H5, between \SI{-38}{dBc} and \SI{-26}{dBc}, are present already
        at \SI{1}{kHz} and come from the weight errors of the DAC.
\end{itemize}}{%
\subsection{Dynamic performance}
The transient simulation of the sine records (the same frequencies as the measurement, with the ideal
differential input and with the inputs of the front-end model of Section~\ref{sec:board}) is still
running; this subsection is completed when it finishes.}

\subsection{Consistency}
\begin{itemize}
  \item \textbf{Consistent}: the converter works as designed at the board conditions: monotonic, no
        missing codes and the LSB within a few percent of the ideal, as in the simulation.
\ifdefined\PexBitEA
  \item \textbf{Explained by the layout}: the bad static linearity. The schematic has binary weights
        within \SI{0.1}{\percent}, the chip errors of \SIrange{-8}{7}{\percent} that cost about one bit of
        ENOB. The post-layout simulation of the taped-out layout has the same sign of error on all five
        bits, \PexInlRms{} of the \InlRms{} LSB rms of measured INL, and about half of the lost ENOB.
        The cause is the wiring capacitance of the DAC plates to the comparator inputs, which differs
        from bit to bit (bridges of bits 1 and 3, MSB plate next to its own top plate) and between the
        two halves of the array. It is a property of the layout, not of the process or of the board.
  \item \textbf{Not yet explained}: the remaining difference (bit 0 less reduced, bit 2 larger and
        bit 3 smaller than post-layout, about \PexCmpInlRms{} LSB rms of INL). Its size is what random
        mismatch of the unit capacitors (not simulated), the wiring resistance (not extracted) or the
        undecoupled references of the board can give; a ramp with the references decoupled at the chip
        would separate the board from the chip.
\else
  \item \textbf{Not consistent}: the static linearity. The schematic has binary weights within
        \SI{0.1}{\percent}; the chip has errors of \SIrange{-8}{7}{\percent} that cost about one bit of
        ENOB. The simulation does not include the layout parasitics of the DAC, the impedance of the
        references on the board or the mismatch of the unit capacitors, and one of these is the cause.
        The post-layout simulation and a ramp measurement with the references decoupled at the chip
        (Section~\ref{sec:board}) will tell which.
\fi
  \item \textbf{Not attributable to the ADC}: the loss of ENOB above \SI{100}{kHz} and the rise of H2,
        which the front end of the board explains. The simulation with the front-end model reproduces
        the loss of ENOB; the H2 needs the nonlinearity of the op-amps on top
        (Section~\ref{sec:board}).
\end{itemize}
